\begin{tikzpicture}[
  font=\small,
  identity/.style={draw=paperblueedge, fill=paperblue, rounded corners=2pt,
    minimum width=2.65cm, minimum height=1.12cm, align=center,
    text=paperink, font=\small\bfseries},
  runtime/.style={draw=paperink, fill=papergray, rounded corners=2pt,
    minimum width=2.65cm, minimum height=1.12cm, align=center,
    text=paperink, font=\small\bfseries},
  resolve/.style={-{Latex[length=2.2mm]}, very thick, draw=paperink},
  local/.style={-{Latex[length=2.2mm]}, thick, dashed, draw=paperblueedge},
  kind/.style={font=\scriptsize, text=paperblueedge, align=center}
]
  \node[identity] (catalog) {Service Model};
  \node[identity, right=0.52cm of catalog] (manifest) {모델 아티팩트\\명세서};
  \node[identity, right=0.52cm of manifest] (grant) {인가 증서};
  \node[runtime, right=0.52cm of grant] (peer) {추론 피어};
  \node[runtime, right=0.52cm of peer] (engine) {서빙 엔진};

  \node[kind, below=0.16cm of catalog] {제품 식별자\\사용자에게 표시};
  \node[kind, below=0.16cm of manifest] {아티팩트 식별자\\서명된 구성};
  \node[kind, below=0.16cm of grant] {제한적 권한\\대상과 범위};
  \node[kind, below=0.16cm of peer] {배포 식별자\\라우팅 대상};
  \node[kind, below=0.16cm of engine] {로컬 어댑터\\응용에 비공개};

  \draw[resolve] (catalog) -- (manifest);
  \draw[resolve] (manifest) -- (grant);
  \draw[resolve] (grant) -- (peer);
  \draw[local] (peer) -- (engine);
\end{tikzpicture}
